% BEGIN TEMPLATE
\documentclass{article}
\usepackage{graphicx}
<<<<<<< HEAD
\usepackage{hyperref} 
\usepackage{xcolor}
\usepackage{nameref}
\usepackage{listings}
\graphicspath{ {../../images/} }
\bibliographystyle{acm}
% CHANGE THESE
\newcommand{\courseListing}{CSCI 8110-001}
\newcommand{\courseName}{Advanced Machine Learning Applications}
\newcommand{\assignmentTitle}{Homework Assignment \#1}
\newcommand{\assignmentSubtitle}{Convolutional Autoencoders}

\hypersetup{
    colorlinks,
    linkcolor={red!50!black},
    citecolor={blue!50!black},
    urlcolor={blue!80!black}
}
\urlstyle{same}
\definecolor{codegreen}{rgb}{0,0.6,0}
\definecolor{codegray}{rgb}{0.5,0.5,0.5}
\definecolor{codepurple}{rgb}{0.58,0,0.82}
\lstdefinestyle{mystyle}{
    commentstyle=\color{codegreen},
    keywordstyle=\color{magenta},
    numberstyle=\tiny\color{codegray},
    stringstyle=\color{codepurple},
    basicstyle=\ttfamily\footnotesize,
    breakatwhitespace=false,         
    breaklines=true,                 
    captionpos=b,                    
    keepspaces=true,                 
    numbers=left,                    
    numbersep=5pt,                  
    showspaces=false,                
    showstringspaces=false,
    showtabs=false,                  
    tabsize=2
}

\lstset{style=mystyle}

\begin{document}
  \begin{center}
  \includegraphics[scale=0.15]{UNO-Logo-Color.png}
  \\[0.3in]
  \textbf{\courseListing{}}\\
  \courseName{}
  \\[0.75in]
  \textbf{\assignmentTitle{}}\\
  \assignmentSubtitle{}
  \\[0.75in]
  \textbf{Patrick Davlin}
  \\[0.75in]
  \textbf{Computer Science Department}\\
  \textbf{Peter Kiewit Institute}\\
  \textbf{University of Nebraska}
  \\[0.75in]
  \textbf{Fall 2020}
  \\[0.3in]
  \includegraphics[scale=0.075]{UNO-Icon-Color.png}
  \newpage
\end{center}
  \graphicspath{{./images/}}
% END TEMPLATE
\section{Introduction}

\section{Project Setup}
\par All work in this project was developed and executed using the Paperspace Gradient service.
Gradient allows for rapid setup of TensorFlow notebooks on GPU-powered machines, which was ideal for this first project given that optimizing a personal computer for running machine learning models proved difficult at the outset.
The downside, however, is that without a more expensive subscription, Gradient only allows for six hours' worth of compute time at once (further discussed in \nameref{procresults}).
This limits the amount of time that computations can run, and Gradient is strict about cutting off at that point--more than one computation for this project was cut off due to poor calculation of the time a task would take to complete.

\par Once the Gradient instance is up and running, the instance is a reasonably straightforward Jupyter Notebook running TensorFlow 2.0.
Blocks of code can be executed interactively and accessed over the web.
A static version of the notebook can be accessed by clicking on \href{https://console.paperspace.com/te7vzjiu3/notebook/prz0iko1d}{this link}.

\section{Process \& Results} \label{procresults}
\par This first assignment proved difficult for two reasons. 
First, with minimal experience in setting up Jupyter notebooks or working with TensorFlow, it was challenging to set up and successfully run code out of the gate in the Paperspace instance.
Some time was spent off the bat learning about some of the conventions of working with Jupyter notebook files to ensure that proper processes were being followed.

\par The second, more project-specific, issue was that adapting the code from class to accept an input size appropriate to the CIFAR-10 dataset was not immediately intuitive.
This was a byproduct of inexperience with TensorFlow and, initially, a lack of understanding of how 

\section{Conclusions}
\newpage
\section{Appendices}
\subsection{Complete Code Listing}

\lstinputlisting[language=Python]{HW1_code.py}

\begin{thebibliography}{9}
  \bibitem{GoAiWired} 
  Cade Metz. 2017.
  \textit{The Sadness and Beauty of Watching Google's AI Play Go}. 
  Retrieved August 27, 2020 from \url{https://www.wired.com/2016/03/sadness-beauty-watching-googles-ai-play-go/}.
  
  \end{thebibliography}

=======
\graphicspath{ {../../images/} }
% CHANGE THESE
\newcommand{\courseListing}{CSCI 8110-001}
\newcommand{\courseName}{Advanced Machine Learning Applications}
\newcommand{\assignmentTitle}{Assignment \#1}
\newcommand{\assignmentSubtitle}{Convolutional Autoencoders}

\begin{document}
  \begin{center}
  \includegraphics[scale=0.15]{UNO-Logo-Color.png}
  \\[0.3in]
  \textbf{\courseListing{}}\\
  \courseName{}
  \\[0.75in]
  \textbf{\assignmentTitle{}}\\
  \assignmentSubtitle{}
  \\[0.75in]
  \textbf{Patrick Davlin}
  \\[0.75in]
  \textbf{Computer Science Department}\\
  \textbf{Peter Kiewit Institute}\\
  \textbf{University of Nebraska}
  \\[0.75in]
  \textbf{Fall 2020}
  \\[0.3in]
  \includegraphics[scale=0.075]{UNO-Icon-Color.png}
  \newpage
\end{center}
  \graphicspath{{./hw1-images/}}
% END TEMPLATE
\section{Introduction}
\textit{[placeholder]}
\section{Project Setup}
All work for this assignment was developed and executed using the Paperspace Gradient online service.
Gradient allows for rapid setup of TensorFlow environments on GPU-powered machines, which was ideal for this first project given that optimizing a personal computer for running machine learning models proved to be difficult for a first attempt at executing the concepts covered in lecture. The downside, however, is that 
\\[0.5in]
\textbf{TODO}: add citation formatting
\\
(1) needs to be https://www.wired.com/2016/03/sadness-beauty-watching-googles-ai-play-go/
\\
(2) needs to be https://towardsdatascience.com/the-truth-behind-facebook-ai-inventing-a-new-language-37c5d680e5a7
>>>>>>> 3ef1ff21c6404548dc895cc789427cebe01e3333
\end{document}